% Physical unit 007; accepted Session 02.
\GordanUnitStart{007}{printed page 2}
\GordanTopFolio{2}
riae hujus circuli in sphaera erectarum, unaquaeque effingit cum aequatore angulum, ceteris aequalem, qui \gq{latitudo reducta} paralleli circuli sphaeroidis appellatur et signo \gq{$l$} significatur. ---

\begin{center}\S\ 6.\end{center}
Circulus enim omnis sphaeroidis, cum et latitudine sua \gq{$B$} et latitudine \gq{$l$} determinetur, appellatur circulus \gq{$B$}, sive circulus \gq{$l$}. Radius circuli \gq{$l$} $=a\cos l$.

\begin{center}\S\ 7.\end{center}
Altera ex latitudine facile determinatur altera. Siquidem inter angulum $l$ et angulum $B$ ratio intercedit haec:
\[
  a\,\tg l=b\,\tg B
\]
vel, cum excentricitas ellipseos meridiani, scilicet $\sqrt{1-\frac{b^2}{a^2}}$, \gq{$e$} significatur, ratio haec:
\[
  \tg l=\sqrt{1-e^2}\,\tg B.
\]
Quae etsi est simplicissima, latitudines $B$ et $l$ alteram ad alteram reducendi, tamen saepenumero aliis rationibus utemur, quas omnes ad superiorem referri notum est.

Quarum rationum summi momenti sunt hae:
\[
\begin{aligned}
\sin l&={b\over a}{\sin B\over\sqrt{1-e^2\sin^2 B}},
&\qquad
\sin B&={\sin l\over\sqrt{1-e^2\cos^2 l}},\\
\cos l&={\cos B\over\sqrt{1-e^2\sin^2 B}},
&
\cos B&={b\over a}{\cos l\over\sqrt{1-e^2\cos^2 l}},\\
{b^2\over a^2}&=1-e^2=(1-e^2\sin^2 B)(1-e^2\cos^2 l).
\end{aligned}
\]
denique posito:
\[
  {a-b\over a+b}={1-\sqrt{1-e^2}\over 1+\sqrt{1-e^2}}=\varepsilon,
\]
\[
  B=l+\varepsilon\sin 2l+{\varepsilon^2\over 2}\sin 4l+{\varepsilon^3\over 3}\sin 6l.
\]

\begin{center}\S\ 8.\end{center}
Quoniam angulis $l$ et $\psi$ quodque sphaeroidis punctum, ut supra diximus, determinatur, illos coordinatarum loco haberi licet. Ut enim cujuslibet puncti situs in plano duabus coordinatis definitur, et plana curva ratione inter has duas coordinatas intercedente, determinatur, ratio inter $l$ et $\psi$ intercedens exprimat curvam in sphaeroide. Ut porro in discussione planarum curvarum extensio et inclinatio in quoque puncto in axem
